%\documentclass[draft]{beamer}
\documentclass[xcolor=pdftex,svgnames,table]{beamer}
%\usepackage[table]{xcolor}
\usepackage{latexsym,amssymb,amsmath,stmaryrd}
\usepackage[utf8]{inputenc}
\usepackage[T1]{fontenc}
\usepackage[francais]{babel}
\usepackage{pgf}
\usepackage{tikz}
\usetikzlibrary{arrows,patterns,plotmarks,shapes,snakes,er,3d,automata,%
backgrounds,topaths,trees,petri,mindmap}
\usepackage{pgfbaseimage}
\usepackage{ifpdf}
\ifpdf
\usepackage{graphicx}
\else
%\usepackage[dvips]{graphicx}
\usepackage{pstricks,pst-tree,pst-node}
\newpsobject{showgrid}{psgrid}{subgriddiv=1,griddots=10,gridlabels=6pt}
\fi
\usepackage{moreverb}
\usepackage[normalem]{ulem}
\usepackage{version}
\usepackage{url}
\usepackage{multicol}
\usepackage{listings}
\lstset{
  language=[ANSI]C, 
%  gobble=2, 
%  escapeinside="", 
  basicstyle=\ttfamily, 
%  directivestyle=\color{Fuchsia},
  identifierstyle = \color{DarkOrange},
  keywordstyle=\color{DarkGreen}, 
  commentstyle=\color{red}, 
%  numbers=left, 
%  numbersep=5pt, 
%  numberstyle=\scriptsize, 
%  moredelim=[is][\color{gray}\itshape]{/*}{*/}, 
%  moredelim=[is][\alert]{/+}{+/}, 
%  morecomment=[is]{/=}{=/}, 
%  morecomment=[in]{comment=}{,} 
%  fancyvrb=true 
} 

\newcommand{\binaire}[1]{\ensuremath{\underline{#1}}}
\newcommand{\C}[1]{\texttt{#1}}
\newcommand{\bbbn}{\ensuremath{\mathbb{N}}}
\input{communs/cc_beamer}
\usetheme{classic}
\newcommand{\nowrite}{\put(10,-4){\includegraphics[scale=.05]{creative_commons/nopencil}}}
\newcommand{\youwrite}{\put(10,-4){\includegraphics[scale=.05]{creative_commons/pencil}}}
\newcommand{\aemporter}{\put(10,-6){\includegraphics[scale=.05]{creative_commons/szymonraj_Shopping_bag}}}


%%% Titre -- Cours 5
\title{Algorithmique et programmation -- Cours 3. Boucle while,
expressions booléennes. Algorithmes élémentaires.}

\author{Pierre Boudes}
\date{\today}

\begin{document}

%% Page de titre et licence CC.
\begin{frame}
	\titlepage
	\vfill
	\begin{center}
		\CcGroupByNcSa{0.83}{0.95ex}\\[2.5ex]
		{\tiny\CcNote{\CcLongnameByNcSa}}
		\vspace*{-2.5ex}
	\end{center}
\end{frame}

\frame[label=plan]{\tableofcontents}

\section{printf/scanf (1)}
\begin{frame}[fragile]
  \frametitle{printf/scanf (1) \nowrite}
\begin{itemize}
\item  Pour afficher un texte à l'écran, nous utilisons la fonction \alert{printf}
  (\emph{print formatted}). 
\item Chaque \verb+%+ 
dans le texte à afficher est substitué par la valeur formatée d'un
\alert{paramètre supplémentaire} de la fonction. Le caractère suivant le
symbole \verb+%+ 
signale le format à utiliser.  Un \verb+%d+
met une valeur au format \alert{entier décimal}.
\pause
\item Exemples :
  \begin{itemize}
\item
    \verb+printf("Bonjour\n")+ 
affiche \C{Bonjour} et un saut de
    ligne\pause
\item
    \verb+printf("i vaut %d\n", i)+ 
affiche \C{i vaut} suivi de la valeur décimale de i (et d'un saut de
ligne)\pause
\item
    \verb+printf("(%d, %d)\n", 31, -4)+ 
affiche \C{(31, -4)} et un saut de ligne.
 \end{itemize}
\pause
\item Réciproquement pour faire entrer dans le programme une donnée saisie par
  l'utilisateur, nous utiliserons \alert{\C{scanf}}.
\item Exemple : \verb+scanf("%d", &x)+
\end{itemize}
\end{frame}



%\subsection{Rappels sur l'instruction de contrôle if}
\begin{frame}[fragile]
  \frametitle{Rappels sur l'instruction de contrôle if \nowrite}

  \structure{Syntaxe :} \C{if (\emph{condition}) \{\emph{\colorbox{green!20}{bloc1}}\} else 
    \{\emph{\colorbox{red!20}{bloc2}}\}}.
  \medskip

  % \pause
  \begin{columns}
    \begin{column}[t]{5cm}
      \structure{Code source}
\begin{verbatim}
/* avant */
if (age < 18)
{
  permis = 0;
}
else
{
  permis = 1;
}
/* après */
\end{verbatim}
    \end{column}
    \pause
    \begin{column}[t]{4cm}
      \structure{Schéma de traduction}\\[.3cm]
      \tikzstyle{amil}=[rectangle, draw,  fill=blue!20, text badly ragged, text width=3.7cm]

      \tikzstyle{line} = [draw, thick, -latex']
      \begin{tikzpicture}[node distance=.7cm]
        \node [amil] (avant) {avant};
        \node [amil,below of=avant] (eval) {évalue la condition};
        \node [amil,below of =eval,fill=orange!20] (selection) {saute si elle est fausse};
        % \uncover<4->{%
        \node [amil,below of =selection,fill=green!20] (bloc1) {bloc1};
        \node [amil,below of =bloc1,fill=orange!20] (saut1) {saute après};
        % }
        \node [amil,below of =saut1,fill=red!20](bloc2) {bloc2};
        \node [amil,below of =bloc2] (apres) {après};
        \path [line] (selection) -| ++(2.6,0) |-  (bloc2);
        % \uncover<4->{
        \path [line] (saut1) -| ++(2.4,0) |-  (apres);
        % }
      \end{tikzpicture}
    \end{column}
  \end{columns}
\end{frame}

\subsection{L'instruction de contrôle for}
\begin{frame}[fragile]
  \frametitle{Rappels sur l'instruction de contrôle for\nowrite}

  \structure{Syntaxe :}\\
  \C{for (\emph{instruct1}; \emph{condition}; \emph{\colorbox{yellow!20}{instruct2}}) 
    \{\emph{\colorbox{green!20}{bloc}}\}}.
  \medskip

  \pause
  \begin{columns}
    \begin{column}[t]{6cm}
      \structure{Code source}
\begin{verbatim}
/* avant */
for (i = 0; i < 5; i = i + 1)
{
  printf("%d\n", i);
  ...
}
/* après */
\end{verbatim}
      La variable \C{i} est appelée \alert{variable de boucle}, elle doit être préalablement déclarée comme toute autre variable.

    \end{column}
    \uncover<3->{
      \begin{column}[t]{4cm}
        \structure{Schéma de traduction}\\[.3cm]
        \tikzstyle{amil}=[rectangle, draw,  fill=blue!20, text badly ragged, text width=3.7cm]
        \tikzstyle{line} = [draw, thick, -latex']
        \begin{tikzpicture}[node distance=.7cm]
          \node [amil] (avant) {avant};
          \node [amil,below of=avant] (instruction1) {instruction 1};
          \node [amil,below of =instruction1,fill=orange!20] (saut1) {saute à la condition};
          \node [amil,below of =saut1,fill=green!20] (bloc) {bloc};
          \uncover<4->{\node [amil,below of =bloc,fill=yellow!20] (instruction2) {instruction 2};}
          \node [amil,below of =instruction2](eval) {évalue la condition};
          \node [amil,below of =eval,fill=orange!20](sautcond) {saute si elle est vraie};
          \node [amil,below of =sautcond] (apres) {après};
          \path [line] (saut1) -| ++(2.6,0) |-  (eval);
          \path [line] (sautcond) -| ++(2.4,0) |-  (bloc);
        \end{tikzpicture}
      \end{column}
    }
  \end{columns}
\end{frame}


\section[while]{L'instruction de contrôle while}
\subsection{Syntaxe et sémantique}
\begin{frame}[fragile]
  \frametitle{L'instruction de contrôle while\youwrite}

\structure{Syntaxe :}\\
 \C{while (\emph{condition}) 
\{\emph{\colorbox{green!20}{bloc}}\}}.
\medskip

\pause
  \begin{columns}
    \begin{column}[t]{6cm}
\structure{Code source}
\begin{verbatim}
/* avant */
while ( ... )
{
 ...
}
/* après */
\end{verbatim}
\uncover<4->{Pour assurer la terminaison, le bloc du while doit
  modifier la condition.}
  \end{column}
\uncover<3->{
   \begin{column}[t]{4cm}
\structure{Schéma de traduction}\\[.3cm]
\tikzstyle{amil}=[rectangle, draw,  fill=blue!20, text badly ragged, text width=3.7cm]
\tikzstyle{line} = [draw, thick, -latex']
     \begin{tikzpicture}[node distance=.7cm]
       \node [amil] (avant) {avant};
      \node [amil,below of =avant,fill=orange!20] (saut1) {saut};
       \node [amil,below of =saut1,fill=green!20] (bloc) {bloc};
      \node [amil,below of =bloc](eval) {évalue la condition};
       \node [amil,below of =eval,fill=orange!20](sautcond) {saute si elle est vraie};
       \node [amil,below of =sautcond] (apres) {après};
       \path [line] (saut1) -| ++(2.6,0) |-  (eval);
       \path [line] (sautcond) -| ++(2.4,0) |-  (bloc);
     \end{tikzpicture}
   \end{column}
}
 \end{columns}
\end{frame}

%\subsection{Trace}
\begin{frame}[fragile]
  \frametitle{Trace}
\begin{columns}
  \begin{column}[t]{3.7cm}
 \begin{lstlisting}[numbers=left,basicstyle=\ttfamily\scriptsize]
int main()
{
    /* Declaration et initialisation des variables */
    int x;
    int nb = 1; /* nombre de chiffres */

    printf("Entrer un nombre positif ");
    scanf("%d", &x);
    
    while (x > 9) /* tant que x a plus d'un chiffre */
    {
        x = x / 10; /* enlever un chiffre a x */
        nb = nb + 1; /* compter un chiffre de plus */
    }
    printf("ce nombre a %d chiffres decimaux\n", nb);
    
    /* Valeur fonction */
    return EXIT_SUCCESS;
}
\end{lstlisting}
\vspace{.4cm}
  \end{column}
\begin{column}[t]{4cm}
\scriptsize
\pause
L'utilisateur saisit 6071.

\rowcolors[\hline]{2}{green!25}{yellow!50} %\arrayrulecolor{green!75!gray}
\begin{tabular}{|r|c|c|l|}
\hline
    ligne & \C{x} & \C{nb} & sortie écran \pause \\ \hline
  initialis.& ? &  1 & \pause \\
7 &  &  & \C{Entrer\ldots}\pause \\
8 &  6071 &  & \pause\\
12 &  607 &  & \pause \\
13 &   & 2 & \pause\\
12 &  60 &  & \pause\\
13 &   & 3 & \pause \\
12 &  6 &  & \pause\\
13 &   & 4 &  \pause \\
15 & & & \C{ce nombre a 4...}\pause \\ 
18 & \multicolumn{3}{c}{Renvoie \C{EXIT\_SUCCESS}}
\end{tabular}
 \end{column}
\end{columns}
\end{frame}

\subsection{For ou while ?}
\begin{frame}
  \frametitle{For ou while ?}

  \begin{itemize}
  \item Un \C{for} peut toujours être simulé par un \C{while} et le
    code machine sera identique. Il suffit d'introduire un
    \alert{compteur de boucle} (la variable de boucle du
    \C{for}).\pause
\item Par convention, les programmeurs préfèrent utiliser un \C{for} lorsque
  le nombre d'itérations est connu à l'avance. Par exemple, pour
  parcourir un ensemble de cas. Dans le cas contraire, les
  programmeurs utilisent un \C{while}. Par exemple, pour chercher un
  cas particulier.\pause
\item Maintenant que nous avons le \C{while}, il est possible qu'un
  programme ne termine jamais (\C{Ctrl-C}).
\item \emph{L'arrêt des programmes est indécidable}.
  \end{itemize}
\end{frame}

\section{Expressions booléennes}

\subsection{Syntaxe}
\begin{frame}
  \frametitle{Expressions booléennes}
Les \emph{conditions} employées dans les structures de contrôle
(\C{if}, \C{for} ou \C{while}) sont des \alert{expressions
booléennes}, pouvant être \emph{Vrai}, \emph{Faux}  ou :\pause
  \begin{itemize}
   \item des inégalités entre expressions arithmétiques
      \begin{align*}
        \text{\emph{inégalité}} &\mathrel{:=} e_1 \mathrel{\C{<}} e_2
        \mid e_1 \mathrel{\C{>}} e_2 \mid e_1 \mathrel{\C{!=}}  e_2 \\
&\mid e_1 \mathrel{\C{<=}}  e_2  \mid e_1 \mathrel{\C{>=}} e_2  \mid e_1 \mathrel{\C{==}} e_2
      \end{align*}
\pause
      \item ou des combinaisons logiques d'expressions booléennes :
        \begin{align*}
          \text{\emph{condition}} &\mathrel{:=}  (\text{\emph{condition}}) \mathrel{\C{\&\&}}
          (\text{\emph{condition}}) \tag{et} \\
&\mid (\text{\emph{condition}}) \mathrel{\C{||}}
          (\text{\emph{condition}})\tag{ou} \\
&\mid {\C{!}}(\text{\emph{condition}}) \tag{non}\\
&\mid \text{\emph{Vrai}} \mid \text{\emph{Faux}} \mid
\text{\emph{inégalité}}\tag{cas de base}
        \end{align*}
  \end{itemize}
\end{frame}


\subsection{Constantes}
\begin{frame}[fragile]
  \frametitle{Constantes booléennes}
  \begin{itemize}
    \item
 Certains langages possédent un type booléen (admettant deux
  valeurs \emph{true} et \emph{false}) pour les expressions
  booléennes.\pause

\item En langage C, les
  expressions booléennes sont de type entier (\C{int}), l'entier \emph{zéro}
  joue le rôle du Faux, l'entier \emph{un} joue le rôle du Vrai et tout
  entier différent de zéro est évalué à vrai.\pause

\item On se donne deux
  constantes symboliques :
\begin{lstlisting}[escapechar={\%},basicstyle=\scriptsize\ttfamily]
/* Declaration des constantes et types utilisateur */
#define TRUE 1
#define FALSE 0

int main()
{
    int continuer = TRUE; /* faut-il continuer ?*/

    while (continuer)
    {
        ...
\end{lstlisting}

  \end{itemize}
\end{frame}


\section{Algorithmes élémentaires}
\newcounter{savedenumi}

\begin{frame}
  \frametitle{Algorithmes : quels outils pour quels problèmes}
\pause
\begin{columns}
 \begin{column}[t]{0.59\linewidth}
  \begin{enumerate}
    \item  Traiter des cas spécifiques 
      \begin{itemize}
      \item \alert{\C{if else}} (différencier)
      \item \alert{\C{\#define}} constantes symboliques (nommer)
      \item arbre de décision (organiser)
     \end{itemize}\pause
    \item Parcourir/générer des cas\pause
      \begin{itemize}
      \item \alert{boucle \C{for}} (rarement while)
      \item tableaux \emph{(plus tard)}
      \end{itemize}\pause
   \item Composer des cas
      \begin{itemize}
        \item boucles (parcourir/générer)
        \pause
      \item \alert{accumulateur} (à initialiser)
      \end{itemize}\pause
   \item Sélectionner des cas
      \begin{itemize}
        \item boucles (parcourir/générer)
        \item \alert{\C{if}}  (sélectionner/traiter)
      \end{itemize}
\setcounter{savedenumi}{\value{enumi}}
\end{enumerate}
\end{column}\hspace{-1cm}
\pause
\begin{column}[t]{0.59\linewidth}
\begin{enumerate}
    \item[3'.] Dénombrer des cas
      \begin{itemize}
        \item boucles (parcourir/générer)
          \pause
        \item \alert{compteur} (à initialiser à 0)
      \end{itemize}\pause
\setcounter{enumi}{\value{savedenumi}}
    \item Rechercher un cas\pause
      \begin{itemize}
        \item boucle \alert{\C{while}}, conditions booléennes, \C{if}
      \end{itemize}\pause
    \item[2'.] Parcourir/générer : une \emph{ligne} mais aussi une \emph{surface}, un \emph{volume}\dots\pause
      \begin{itemize}
        \item imbriquer les boucles (var. $\neq$)
      \end{itemize}\pause
    \item Boucle événementielle
  \begin{itemize}
\item boucle \alert{while}
  \end{itemize}\pause
    \item Attente active 
  \begin{itemize}
\item boucle \alert{while}
  \end{itemize}
\end{enumerate}
\end{column}
\end{columns}
\end{frame}


\subsection{Recherche d'un diviseur (test de primalité)}

\begin{frame}[fragile]
\frametitle{Recherche d'un diviseur (test de primalité)}

\begin{block}{Exemple}
  Le programme cherche un entier supérieur à $1$ qui divise $n$, s'il n'en trouve pas
 de plus petit que $n$, alors $n$ est premier (ou bien $n = 1$).\nowrite
\end{block}
\lstinputlisting[basicstyle=\scriptsize\ttfamily,linerange={31-50},numbers=left,firstnumber=31]{../demo/premier.c}
\end{frame}


\subsection{Autres usages du while : boucle événementielle et attente active}

\begin{frame}[fragile]
\frametitle{Boucle principale en programmation événementielle}

Le programme lit les événements
(les actions de l'utilisateur) et les traite, continuellement, dans
une boucle.

\begin{block}{Remarque}
 \C{scanf} est un \alert{appel système bloquant} (processus en attente).
\end{block}
\begin{block}{Exemple}
Jeu de calcul mental (additionner deux nombres au hasard)\nowrite
\end{block}
\lstinputlisting[basicstyle=\scriptsize\ttfamily,linerange={16-16},numbers=left,firstnumber=16]{../demo/addition.c}
\onslide*<2->{\lstinputlisting[basicstyle=\scriptsize\ttfamily,linerange={22-23},numbers=left,firstnumber=22]{../demo/addition.c}}
\lstinputlisting[basicstyle=\scriptsize\ttfamily,linerange={31-32},numbers=left,firstnumber=31]{../demo/addition.c}
\onslide*<2->{\lstinputlisting[basicstyle=\scriptsize\ttfamily,linerange={34-35},numbers=left,firstnumber=34]{../demo/addition.c}}
\lstinputlisting[basicstyle=\scriptsize\ttfamily,linerange={53-59},numbers=left,firstnumber=53]{../demo/addition.c}
\end{frame}

\begin{frame}[fragile]
  \frametitle{Boucle d'attente active}
La boucle sert à attendre qu'une condition externe devienne vraie, le
programme teste continuellement cette condition. 
\pause
\begin{block}{Exemple}
  \emph{Tic tac boum}, est-ce que 10 secondes sont écoulées ?\nowrite
\end{block}
\lstinputlisting[basicstyle=\scriptsize\ttfamily,linerange={19-23},numbers=left,firstnumber=19]{../demo/tictacboum.c}
\lstinputlisting[basicstyle=\scriptsize\ttfamily,linerange={42-44},numbers=left,firstnumber=42]{../demo/tictacboum.c}
\end{frame}

\section{Démos}
\begin{frame}
  \frametitle{Démos et utilisation de codeblocks}
\end{frame}
\end{document}
